
\chapter{Introduction}
\label{sec:intro}


\chapterintro


Cette thèse se place dans le contexte de vérification et validation des
logiciels.
Les principales approches de ce domaine sont présentées en partie~\ref{sec:vv},
nous y présentons notamment les outils que nous utilisons dans nos travaux :
\pathcrawler en partie~\ref{sec:test-structurel} et \framac en
partie~\ref{sec:framac}.
Nous expliquons la problématique de nos travaux et nos motivations en
partie~\ref{sec:pb-motiv}.
Les contributions de cette thèse sont présentées en partie~\ref{sec:contrib} et
le plan de la thèse est annoncé en partie~\ref{sec:plan}.


\section{Vérification et validation de programmes}
\label{sec:vv}


Le domaine de la vérification et de la validation regroupe un ensemble de
techniques du cycle de développement des logiciels qui ont pour objectif de
s'assurer de leur correction et de leur sûreté. Ces deux notions sont apparues
dans les années soixante-dix avec les travaux de \bsc{Dijkstra}
\cite{Dijkstra/75}, \bsc{Floyd} \cite{Floyd/63} et \bsc{Hoare} \cite{Hoare/69}.

La correction d'un logiciel représente le respect de l'implémentation par
rapport aux spécifications. La sûreté d'un logiciel est liée à son absence
d'erreurs à l'exécution.


\subsection{Model-checking}
\label{sec:model-checking}


Le model-checking \cite{Clarke/86} permet de vérifier algorithmiquement si
un modèle donné (le système ou une abstraction de ce système) satisfait une
spécification, formulée en termes de logique temporelle \cite{Clarke/82}.
Un modèle est un ensemble d'états, de propriétés que vérifie chaque état, et de
transitions entre ces états qui décrivent l'évolution du système.

Le model-checking couvre l'ensemble des états du système et des transitions afin
d'analyser toutes les exécutions possibles du système. Sur de grands systèmes,
cette méthode est pénalisée par l'explosion combinatoire du nombre des états (et
la complexité en temps ou en espace qui en résulte). Il est néanmoins possible
de modéliser des algorithmes asynchrones répartis et des automates non
déterministes, comme le fait notamment l'outil \spin \cite{\citespin}.

%% La plateforme \framac intègre \textsc{Aoraï}, un greffon
%% permettant d'annoter automatiquement le code source d'un programme C d'après
%% une formule de logique temporelle linéaire, de sorte que les annotations sont
%% vérifiées si le programme repecte la formule.


\subsection{Analyse statique}
\label{sec:AS}


L'analyse statique \cite{Nielson/99} examine le code source du programme
sans l'exécuter. Elle raisonne sur tous les comportements qui pourraient
survenir lors de l'exécution et permet donc de déduire des propriétés devant
être vérifiées pour toutes ces exécutions, dans le but de prouver la correction
du programme.

En revanche, la vérification de programme étant en général indécidable
\cite{Landi/92}, il est souvent nécessaire d'utiliser des
sur-approximations, ce qui implique que les résultats peuvent être moins précis
que ce que l'on souhaite mais ils sont garantis pour toutes les exécutions.
Ainsi, on peut établir des propriétés de sûreté ({\em safety}), où l’on cherche
des invariants sur les valeurs des variables du programme (une plage de valeurs
par exemple), afin d'exclure certains risques d'erreurs à l'exécution.

Parmi les méthodes statiques sont distinguées : l'interprétation abstraite
(section~\ref{sec:interpretation-abstraite}), l'abstraction à base de
prédicats (section~\ref{sec:abstraction-predicats}) et la preuve de programmes
(section~\ref{sec:preuve}).


\subsubsection{Langages de spécification}
\label{sec:speclang}


\commentGP{TODO}
Java Modeling Language (\jml)
+
ANSI/ISO C Specification Language (\acsl)


\subsubsection{Interprétation abstraite}
\label{sec:interpretation-abstraite}


L'interprétation abstraite \cite{Cousot/92} s'appuie sur les
théories du point-fixe et des domaines pour introduire des sur-approximations
des comportements d'un programme. Elle consiste à abstraire les domaines des
variables par des domaines finis et beaucoup plus petits. Par exemple, le
domaine des entiers pourrait être abstrait par un domaine de trois valeurs :
$(-, 0, +)$. Appliquée à l’analyse de valeurs, elle consiste à calculer à
chaque ligne du code une sur-approximation de l’ensemble des valeurs prises par
chaque variable en cette ligne lors de toutes les exécutions du programme,
permettant ainsi de détecter certaines erreurs comme les divisions par zéro ou
les accès en dehors des bornes des tableaux.

Pour contourner le problème d’indécidabilité, la théorie de l’interprétation
abstraite construit une méthode qui, à la même question, répondra ``oui'',
``non'' ou ``peut-être''. Si la méthode répond ``peut-être'', c’est qu’on n’a pu
prouver ni l’un ni l’autre des deux premiers cas. C’est ce qu’on appelle une
alarme : il est possible qu’une des exécutions du programme produise une erreur
donnée, mais nous n’avons été capable ni de le confirmer ni de l’infirmer.
L’erreur signalée par une alarme peut ne jamais apparaître à l’exécution, dans
ce cas on l’appelle fausse alarme. On ne calcule donc pas la propriété exacte
mais une abstraction de cette propriété, en imposant la contrainte de sûreté
suivante : ``la propriété abstraite calculée ne doit oublier aucune exécution
concrète''. L'abstraction est effectuée à base de prédicats atomiques
définissant des abstractions des domaines des variables.

\polyspace \cite{\citepolyspace} a été le premier outil commercial
utilisant l'interprétation abstraite pour détecter les erreurs à l'exécution
dans les programmes en C, C++ et Ada mais signale beaucoup de fausses alarmes.
L'ENS a développé \astree \cite{\citeastree}, spécifique au langage C et aux
logiciels critiques. \fluctuat \cite{\citefluctuat},
qui mesure précisément les approximations faites à l'exécution d'un programme C.
\framac intègre un greffon d'interprétation abstraite : \Value
\cite{\citevalue}.


\subsubsection{Abstraction à base de prédicats}
\label{sec:abstraction-predicats}


L'abstraction à base de prédicats \cite{Schiller/OOPSLA12} est une
technique permettant de générer automatiquement des abstractions de systèmes au
nombre d'états infini. Pour un programme $P$ au nombre d'états infini, un
ensemble fini de prédicats $E = \{f_1, ..., f_n\}$ est défini, ces prédicats
sont des expressions booléennes sur les variables de $P$ et les constantes du
langage.

Chaque état concret de $P$ est mis en correspondance avec un état abstrait de
l'abstraction de $P$, après évaluation par les prédicats de $E$. Un état
abstrait est un $n$-upplet de valeurs booléennes correspondant à la
satisfaisabilité des $n$ prédicats (au moyen d'un solveur SMT).

Par exemple, si 3 prédicats $f_1, f_2, f_3$ sont définis et que la
satisfaisabilité de ces prédicats à l'état concret $e$ est évaluée
respectivement à $false, true, true$, alors dans l'abstraction générée, l'état
$e$ correspond à l'état abstrait $(\lnot f_1, f_2, f_3)$.

L'abstraction générée comporte un nombre fini d'états (au plus $2^n$) car il n'y
a qu'un nombre fini de prédicats, le model-checking peut donc être appliqué à
cette abstraction. Si une propriété de sûreté est vérifiée dans l'abstraction,
elle l'est également dans le système concret.

Cette technique est notamment utilisée par \slam \cite{\citeslam}.


\subsubsection{Preuve de programmes}
\label{sec:preuve}


La preuve de programmes utilise des fondements mathématiques et logiques
\cite{Hoare/69} pour prouver des propriétés de programmes. Tout d'abord, le
système
est décrit par un ensemble d'axiomes et de règles d'inférence. Puis, le calcul
de la plus faible précondition \cite{Dijkstra/75} est utilisé pour générer des
formules appelées obligations de preuve, qui sont finalement soumises à un
prouveur de théorèmes, qui applique différentes techniques de résolution.

Contrairement au model-checking, la preuve a l'avantage d'être
indépendante de la taille de l'espace des états, et peut donc s'appliquer sur
des systèmes de grande taille. En contre-partie, cette technique requiert une
expertise de l'utilisateur pour adapter le programme à la preuve (en l'annotant
par exemple) et guider le prouveur si nécessaire.

Il existe des prouveurs automatiques tels que \simplify \cite{\citesimplify},
\ergo \cite{\citeergo} et \zthree \cite{\citezthree}; et des
prouveurs interactifs, où la preuve est guidée par l'utilisateur, tels que
\coq \cite{\citecoq}, \isabelle \cite{\citeisabelle}, et \hol \cite{\citehol}.
Certains de ces prouveurs ou assistants de preuve sont intégrés à
d'autres outils tels \boogie \cite{\citeboogie} ou \escjava \cite{\citeescjava}.
\framac intègre le greffon de preuve, \Wp \cite{\citewp} qui traite des
programmes dont le code contient des annotations \acsl \cite{\citeacsl}.


\subsection{Analyse dynamique}
\label{sec:AD}


L’analyse dynamique est basée sur des techniques d’exécution du programme, de
simulation \cite{Whitner/WSC89} d’un modèle ou d'exécution symbolique
\cite{Clarke/76}, regroupées sous le terme générique ``test''.

Les tests peuvent s’appliquer tout au long du cycle de développement d’un
logiciel. Les tests unitaires vérifient le bon fonctionnement des différentes
entités d’un système, indépendamment les unes des autres. Les tests
d'intégration vérifient la bonne communication entre ces entités. Les tests de
validation s'assurent que les fonctionnalités correspondent au besoin de
l’utilisateur final. Enfin, les tests de non-régression vérifient que l'ajout de
nouvelles fonctionnalités ne détériore pas les anciennes fonctionnalités.

En général, les techniques de test ne sont pas exhaustives et n'explorent qu'un
sous-ensemble des chemins d'exécutions du programme, en conséquence, l’absence
d’échecs lors du passage des tests n’est pas une garantie de bon fonctionnement
du système. Néanmoins, selon les critères utilisés pour la génération des
tests, et selon la couverture des chemins d'exécution fournie par les tests, un
système ainsi validé peut acquérir un certain niveau de confiance.

Les méthodes de test peuvent être classées en trois catégories : le test
aléatoire, le test structurel (section~\ref{sec:test-structurel}) et le test
fonctionnel (section~\ref{sec:test-fonctionnel}). Comme son nom l'indique, le
test aléatoire consiste à générer des valeurs d'entrée du programme au hasard et
ne sera pas détaillé dans cette thèse.


\subsubsection{Langages de spécification exécutables}
\label{sec:execspeclang}


\eacsl est un sous-ensemble ``exécutable'' du langage \acsl
implémenté dans \framac. Contrairement à \acsl, chaque
spécification \eacsl est exécutable : elle peut être évaluée à
l'exécution.


\subsubsection{Exécution symbolique pour le test structurel}
\label{sec:exec-sym}


\commentGP{TODO}


\subsubsection{Test structurel}
\label{sec:test-structurel}


Le test structurel, ou test ``boîte blanche'', est une technique de test qui
fonde la détermination des différents cas de test sur une analyse de la
structure du code source du programme étudié. On distingue deux types de tests
structurels : le test orienté flot de contrôle et le test orienté flot de
données.

Le test orienté flot de données cherche à couvrir certaines relations entre la
définition d’une variable et son utilisation, par exemple, on peut souhaiter
couvrir toutes les lectures d'une variable suivant une écriture.

Le test orienté flot de contrôle s’intéresse quant à lui à la structure du
programme : l'ordre dans lequel les instructions sont exécutées. Il se base sur
le graphe de flot de contrôle du programme : un graphe connexe orienté avec un
unique n\oe{}ud d’entrée et un unique n\oe{}ud de sortie, dont les n\oe{}uds
sont les blocs de base du programme et les arcs représentent les branchements
(conditions). Une couverture structurelle de ce graphe est recherchée, selon un
critère qui peut être par exemple ``toutes les instructions'',
``toutes les branches'' (toutes les décisions), ``tous les chemins'' ou
``tous les $k$-chemins''.

L’exécution symbolique dynamique, ou exécution ``concolique'', associe
l’exécution concrète du programme et l’exécution symbolique afin d’explorer les
chemins du programme. L’exécution concrète sert à confirmer que le chemin
parcouru est bien celui pour lequel le cas de test exécuté a été généré.

Plusieurs outils se basent sur l'exécution concolique pour explorer un programme
sous test, dont \smart \cite{\citesmart}, \pex \cite{\citepex},
\sage \cite{\citesage}, \cute \cite{\citecute}, \klee \cite{\citeklee},
\Exe \cite{\citeexe}, \pathcrawler \cite{\citepathcrawler}. Ces outils
utilisent des solveurs de contraintes pour générer des cas de test permettant
d'aboutir à une couverture souhaitée des exécutions du programme par les tests.

\sage \cite{\citesage} utilise une analyse dynamique afin de générer
des tests pour des programmes au format binaire x86. Il combine une exécution
symbolique à du {\em fuzz testing}. Une première exécution avec des entrées
valides permet de récupérer une trace d'exécution du programme. L'exécution
symbolique de cette trace permet de collecter les contraintes du chemin
d'exécution. De nouveaux chemins sont générés par négation des contraintes.

\pathcrawler \cite{\citepathcrawler} est un outil de génération de tests
structurels pour les programmes C, accessible sous la forme d'un service web :
\textsc{PathCrawler Online} \cite{\citepconline}.

Étant donné un programme C sous test $p$ et une pré-condition sur ses entrées,
il génère des cas de test respectant un critère de couverture de test. Le
critère \emph{tous les chemins} impose une couverture de tous les chemins
faisables de $p$. L'exploration exhaustive de tous les chemins étant en pratique
irréalisable sur des programmes réels, le critère \emph{tous les k-chemins} a
été défini, il limite l'exploration aux chemins qui ont au plus $k$ itérations
consécutives de chaque boucle.

\pathcrawler commence par construire une version instrumentée de $p$
permettant de tracer l'exécution de chaque cas de test, puis il génère les
contraintes représentant la sémantique de chaque instruction de $p$. La
prochaine étape est la génération et la résolution de contraintes pour produire
les cas de test pour un ensemble de chemins $\Pi$ satisfaisant le critère de
couverture. La résolution de contraintes s'effectue à l'aide
d'\eclipse \prolog \cite{\citeeclipse}, un environnement de programmation en
logique par contraintes basé sur \prolog.

Étant donné un préfixe de chemin $\pi$, c'est-à-dire un chemin partiel de $p$,
l'idée est de résoudre les contraintes correspondant à l'exécution symbolique
de $p$ en suivant le chemin $\pi$.
 
La méthode de génération de test est composée des étapes suivantes :

\begin{itemize}
\item[$(\mathcal{G}_1)$]
Création d'une variable logique pour chaque entrée.
Prise en compte des contraintes de la pré-condition.
Le préfixe de chemin initial $\pi$ est vide.
Aller à $(\mathcal{G}_2)$.

\item[$(\mathcal{G}_2)$]
Exécuter symboliquement le chemin $\pi$ : ajout des contraintes et
mise à jour de la mémoire en fonction des instructions de $\pi$.
Si certaines contraintes sont insatisfiables, aller à $(\mathcal{G}_5)$.
Sinon, aller à $(\mathcal{G}_3)$.

\item[$(\mathcal{G}_3)$]
Appeler le solveur de contraintes pour générer un cas de test $t$ satisfaisant
les contraintes du chemin courant. Si les contraintes sont insatisfiables, aller
à $(\mathcal{G}_5)$.
Sinon, aller à $(\mathcal{G}_4)$.

\item[$(\mathcal{G}_4)$]
Exécuter le programme avec trace sur le cas de test $t$ généré pour obtenir
le chemin d'exécution, qui doit commencer par $\pi$.
Aller à $(\mathcal{G}_5)$.

\item[$(\mathcal{G}_5)$]
Calculer le prochain chemin partiel $\pi$ à couvrir. Un parcours en profondeur
détermine la dernière décision $d$ pour laquelle il reste une branche à
explorer. S'il n'existe pas une telle décision, l'algorithme s'arrête. Sinon,
$\pi$ est recalculé et contient maintenant le chemin partiel précédent dans
lequel les contraintes correspondant à $d$ ont été niées, et retour à l'étape
$(\mathcal{G}_2)$. Cela nous assure que tous les chemins faisables sont couverts
(en considérant que le solveur de contraintes peut trouver une solution dans un
temps raisonnable) et que seulement le plus court des préfixes infaisables de
chaque chemin infaisable est exploré.
\end{itemize}


\subsubsection{Test fonctionnel}
\label{sec:test-fonctionnel}


Le test fonctionnel, ou test ``boîte noire'', génère des jeux de test en
fonction du comportement attendu du programme : un cas de test sera choisi pour
chaque comportement particulier. Le test fonctionnel est utilisé pour vérifier
la conformité des réactions du logiciel avec les attentes de l'utilisateur, sans
connaissance du code source. Il existe de nombreuses techniques qui se
différencient par la manière de choisir les données de test, parmi lesquelles :

\begin{description}
\item[le test de partition] \hfill \\
les valeurs d’entrées du logiciel sont regroupées en classes d’équivalence, sur
lesquelles le logiciel doit avoir le même comportement ({\em domain splitting}),
une seule valeur aléatoire est choisie dans chaque classe de la partition;
\item[le test aux limites] \hfill \\
les données de test sont choisies aux bornes des domaines de définition des
variables.
\end{description}

\gatel \cite{\citegatel} est un générateur de
tests fonctionnels qui se base sur une représentation symbolique des états du
système : le programme, les invariants et les contraintes décrivant l'objectif
de test sont exprimés dans le langage \lustre \cite{\citelustre}. Cet outil
offre la possibilité de réaliser des partitions de domaines.


\subsection{Monitoring}
\label{sec:monitoring}


\commentGP{TODO}
RAC


\subsection{\framac}
\label{sec:framac}


\framac \cite{\citeframac} est une plate-forme dédiée à l'analyse statique
des programmes C, conjointement dévelopée par INRIA et le CEA LIST. Son
architecture comporte un noyau et un écosystème de
greffons, rendant l’outil extensible. Les greffons peuvent échanger des
informations et utiliser les  services fournis par le noyau, permettant ainsi
une collaboration entre différentes analyses.

\framac est basé sur \cil \cite{\citecil}, une bibliothèque qui normalise des
programmes C (ISO C99) en opérant des modifications syntaxiques : normalisation
des boucles en utilisant la structure \lstinline{while}, unique
\lstinline{return} pour
chaque fonction, etc. \framac étend \cil pour supporter des annotations
dédiées portant sur le code source, exprimées dans le langage \acsl.
\acsl \cite{\citeacsl} est un langage formel de spécification
comportementale \cite{Hatcliff/12}, inspiré de \jml \cite{\citejml}, pouvant
exprimer des propriétés fonctionnelles de programmes C : pré-conditions,
post-conditions, invariants, etc.

En effet, la spécification d'une fonction comprend les pré-conditions requises
(exprimées par une clause \lstinline{requires}) lors de l'appel et les
post-conditions assurées (\lstinline{ensures}) lors du retour. Parmi ces
post-conditions, une clause indique quels sont les emplacements mémoire qui
peuvent être affectés (\lstinline{assigns}) par la fonction.

\lstinputlisting[caption={Exemple de spécification \acsl},label=lst:swap]
   {listings/swap.c}

Considérons par exemple une spécification fournie pour une fonction
\lstinline{swap}
(Listing~\ref{lst:swap}). La première pré-condition établit que les deux
arguments doivent être des pointeurs valides, autrement dit, le déréférencement
de $a$ ou de $b$ ne produira pas d'erreur à l'exécution. La seconde
pré-condition impose que les emplacements mémoire occupés par chacune de ces
variables soient disjoints. En plus de \lstinline{\valid} et
\lstinline{\separated}, \acsl fournit de nombreux prédicats et
fonctions afin de décrire les états de la mémoire. \lstinline{\at(e,l)}
fait référence à la valeur de l'expression \lstinline{e} à l'état de la mémoire
au label \lstinline{l}. \lstinline{Pre} est un label prédéfini qui fait
référence à l'état de la mémoire avant l'exécution de la fonction. Ainsi, les
post-conditions (\lstinline{ensures}) signifient qu'à la fin de la fonction,
\lstinline{*a} aura la valeur que \lstinline{*b} avait au
début de la fonction, et réciproquement.

\acsl offre aussi la possibilité d'écrire des annotations dans le code
source, permettant d'exprimer des propriétés devant être vraies à un point donné
du programme : les assertions (\lstinline{assert}).
Il est également possible d'exprimer des propriétés devant être vraies avant une
boucle et après chaque itération de cette boucle : les invariants de boucle
(\lstinline{loop invariant}).

Les annotations du langage \acsl sont écrites en utilisant la logique
du premier ordre, et il est possible de définir ses propres fonctions et
prédicats.
Les greffons peuvent valider ou invalider les propriétés \acsl et
générer des annotations \acsl, les annotations sont donc un moyen
d'échanger des informations entre les différentes analyses opérées par les
greffons.


\section{Problématique et motivations}
\label{sec:pb-motiv}


\begin{figure}[bt]
  \begin{center}
    \begin{tikzpicture}[>=latex,font=\scriptsize]
      \node[] at (0,0) (Bob) {\includegraphics[scale=0.15]{figures/user_m.png}};
      \node[below right of=Bob] {\textbf{Bob}, dev. engineer};
      \node[draw=greenv,thick,minimum width=6.7cm,rounded corners,
        align=center] (spec) at (2,1.9) {Requirements/Informal Specification};
      \node[above of=Bob,draw,node distance=1cm,thick,
        rounded corners,minimum width=2.65cm] (code) {Code};
      \node[right of=Bob,node distance=4cm] (Alice) {
        \includegraphics[scale=0.15]{figures/user_f.png}};
      \node[below right of=Alice] {\textbf{Alice}, validation engineer};
      \node[above of=Alice,draw,node distance=1cm,thick,
        rounded corners,minimum width=2.65cm] (annot) {
        Formal Specification};
      \node[draw,dotted,thick,fit=(code) (annot)] (code-annot) {};
      \node[right of=code-annot,draw,node distance=6.5cm,thick,
        rounded corners,minimum width=2.65cm,align=center] (wp) {
        Deductive Verification};
      \draw[->,thick] (code-annot.east) to[out=0,in=90] (wp);
      \draw[->,thick,color=orange] (wp.south) to[out=-90,in=0]
      node[above=.2cm]{unknown}(Alice.east);
    \end{tikzpicture}
  \end{center}
\end{figure}


\section{Contributions}
\label{sec:contrib}


\commentGP{TODO}


Liste des contributions :
\begin{itemize}
\item extension d'une méthode de génération de tests structurels de programmes C
  à des programmes C annotés par le sous-ensemble du langage ACSL appelé E-ACSL
\item développement d'une méthode de combinaison Test et Preuve
\item proposition de scénarios types d'utilisation de cette méthode
\item implémentation de la combinaison Test et preuve dans STADY
\item expérimentations évaluant l'impact de la méthode
\item stage E-ACSL
\end{itemize}


\section{Plan de la thèse}
\label{sec:plan}


\commentGP{TODO}
